The incorporation of radiative corrections is also essential for processes in a kinematical regime that is different from the previous discussions.
Recently, it has been gradually recognized that in order to achieve desired precision in experiment, radiative corrections must be incorporated as a part of MC generators (see also Sec.~\ref{sec:Implementations} for more on MC concepts), and subsequently, the demand for such corrections has increased.
Such an awakening resides in the fact that, typically, the QED NLO radiative corrections compete in size with the hadronic effects.
And thus, once hadronic parameters are to be extracted from data, ignoring radiative corrections would lead to meaningless results, the examples of which we will see below.

Historically, the approach to radiative corrections varied.
In some cases, they would be left out entirely.
If an effort was done to include at least some of the relevant corrections, others that were not negligible were ignored.
In some cases, approximate approaches (soft-photon limit or leading logarithms) were used, which often led to unreliable or misleading results.
All the above-mentioned cases introduce artificial discrepancies between theory and experiment, as also discussed in \secref{RoPTcomp}.
In other words, what is then considered to be ``measured observables'' or related ``measured hadronic parameters'' may include an unsubtracted QED part, which can turn out to be quite a significant contribution.

In the following, we will discuss the current status of QED radiative corrections in some sample decays of the lightest pseudoscalar mesons and baryons and their direct application in experiment.


\subsection{Radiative corrections for \texorpdfstring{$\pi^0$}{pi0} decays}

The NLO QED radiative corrections are now worked out for all the relevant neutral-pion decay channels: the $\pi^0$ Dalitz decay $\pi^0\to e^+e^-\gamma$, the $\pi^0$ double-Dalitz decay $\pi^0\to 2(e^+e^-)$, and the $\pi^0$ rare decay $\pi^0\to e^+e^-$.
We will address these individually in the following subsections.


\subsubsection{\texorpdfstring{$\pi^0\to e^+e^-$}{pi0->ee}}

Let us start with the last-mentioned process.
Regarding the branching ratio, this channel is loop- and helicity-suppressed with respect to the radiative decay $\pi^0\to\gamma\gamma$ by eight orders of magnitude.
It has thus been speculated that the $\pi^0\to e^+e^-$ decay might be sensitive to possible effects of new physics.
And it indeed seemed to be the case when the KTeV Collaboration published the results of their analysis on the precise measurement of the branching ratio in 2007~\cite{KTeV:2006pwx}.
The direct subsequent comparison to the Standard Model (SM) prediction was then interpreted as a 3.3\,$\sigma$ discrepancy~\cite{Dorokhov:2007bd}.
Many works followed trying to explain this difference both within (via modeling of the electromagnetic pion transition form factor) and outside the SM (by introducing various models including exotic particles).
However, it is believed that the most important contribution in this regard was done by an explicit calculation of (two-loop) virtual radiative corrections~\cite{Vasko:2011pi} that brought the discrepancy down to the $2\,\sigma$ level~\cite{Husek:2015wta}: Until then, only leading-log approximation was available~\cite{Bergstrom:1982wk,Dorokhov:2008qn} that did not reproduce well accidental cancellations among various diagrams.
New light can be shed on the branching ratio determination and the related remaining $2\,\sigma$ discrepancy: A new measurement is under preparation by the NA62 Collaboration and the results should be known, conservatively estimating it, within a few months, i.e., by the end of 2023.
In this analysis, all the known related radiative corrections~\cite{Vasko:2011pi,Husek:2014tna,Husek:2015sma,Husek:2018qdx} are already included at the MC level.

Looking at the quantity that has been measured by KTeV~\cite{KTeV:2006pwx},
%
\begin{equation}
\frac{B\bigl(\pi^0\to e^+e^-(\gamma),x>0.95\bigl)}
{B\bigl(\pi^0\to e^+e^-\gamma,x>0.232\bigl)}
=1.685(64)(27)\times10^{-4}\,,
\label{eq:BBKTeV}
\end{equation}
%
with $x\equiv m_{e^+e^-}^2/M_\pi^2$, it becomes immediately obvious why there is a need for precise knowledge of radiative corrections in order to compare this value with the predictions provided by theory, which can be written as
%
\begin{equation}
B\bigl(\pi^0\to e^+e^-\bigr)
\approx\bigl(6.21+0.15\widetilde\chi\bigr)\times10^{-8}\,,
\end{equation}
%
with $\widetilde\chi\equiv2\big[\chi^{(\mathrm{r})}(770\,\mathrm{MeV})-\frac52\big]$.
Theoretical predictions and models suggest (being rather conservative) $\chi^{(\mathrm{r})}(770\,\mathrm{MeV})\sim$\;2--3 (see, e.g., Refs.~\cite{Knecht:1999gb,Dorokhov:2007bd,Husek:2015wta,Hoferichter:2021lct}), so $\widetilde\chi$ is expected to be within the interval $(-1,1)$.
The result in Eq.~\ref{eq:BBKTeV} not only depends on the Dalitz-decay branching ratio serving here as the normalization channel, but the numerator allows for final states including soft photons with their energies limited by the requirement on the minimum of the normalized electron--positron invariant mass squared $x$.
Eq.~\ref{eq:BBKTeV} thus needs to be further processed and it was done so in the KTeV analysis via a series of steps in order to provide a ``no-radiation'' value to be compared with theory.
In other words, experimental results on the branching ratios of such processes will contain effects that need to be subtracted in some way, so they can be directly compared with theoretical predictions for branching ratios with no final-state radiation.
Such a subtraction of radiative corrections can happen at different stages of the experimental analysis, but it seems that the earlier stage this happens, the better:
Once the events accounting for photons in the final state are properly generated already at the MC level, comparing the resulting spectra with data helps to improve the quality of the analysis outcome.


\subsubsection{\texorpdfstring{$\pi^0\to e^+e^-\gamma$}{pi0->eeg}}

We have already seen that in the KTeV rare-decay search described above, the Dalitz decay was used as the normalization channel.
It is thus clear that it is essential to have the complete set of NLO QED corrections available also in this case.
It then comes as a bonus that this process is also used for normalizing purposes in measurements in the kaon sector~\cite{Husek:2018qdx}.
Furthermore, investigating the Dalitz decay itself leads to information about the singly off-shell ($\pi^0\gamma\gamma^*$) electromagnetic transition form factor (TFF) in the time-like region, i.e., direct access to the TFF slope without the need for model-dependent extrapolation from the space-like region, as used, e.g., in Ref.~\cite{CELLO:1990klc}.

The radiative corrections for $\pi^0\to e^+e^-\gamma$ were first estimated in terms of the total correction to the integral decay width~\cite{Joseph:1960zz}, and only a decade later, the first corrections to the differential decay width appeared, which used the soft-photon approximation~\cite{Lautrup:1971ew}.
What can now be considered as a classical work on the Dalitz-decay radiative corrections came a year after~\cite{Mikaelian:1972yg}.
The hard-photon corrections were included and the results were presented in terms of a table giving sample correction values throughout the Dalitz plot.
One contribution was, however, missing: the two-photon-exchange contribution, or also known as the one-photon-irreducible (1$\gamma$IR) contribution at one loop, which diagrammatically coincides with the bremsstrahlung contribution to the $\pi^0\to e^+e^-$~\cite{Husek:2014tna}; see Fig.~\ref{fig:1gIR}.
%
\begin{figure}[ht]
    \centering
    \begin{subfigure}[t]{0.48\columnwidth}
        \centering
        \includegraphics[width=\textwidth]{Mesons/fig/T2.eps}
    \end{subfigure}
    %
    \begin{subfigure}[t]{0.48\columnwidth}
        \centering
        \includegraphics[width=\textwidth]{Mesons/fig/T4.eps}
    \end{subfigure}\\[3mm]
    \caption{One-photon-irreducible Feynman diagrams contributing to neutral-pion Dalitz-decay radiative corrections.}
    \label{fig:1gIR}
\end{figure}
%
It turned out, after a vivid discussion in the literature~\cite{Tupper:1983uw,Lambin:1985sb,Tupper:1986yk}, that this contribution was not negligible, and a direct subsequent calculation would show that omitting the 1$\gamma$IR piece would introduce $\sim$\,15\% discrepancy to the measured slope~\cite{Kampf:2005tz,Husek:2015sma}.

The most precise measurement of the slope in the time-like region to date was performed by the NA62 Collaboration with the results $a_\pi=3.68(57)\%$~\cite{NA62:2016zfg}.
It is worth mentioning that the slope of the radiative corrections amounts to circa $-12.5\%$; see the low-$x$ region of Fig.~\ref{fig:Dalitz_RC}.
%
\begin{figure}[ht]
    \centering
    \includegraphics[width=\columnwidth]{Mesons/fig/delta_all.pdf}\\[4mm]
    \caption{
    Inclusive NLO QED radiative corrections for the discussed Dalitz decays $A\to Be^+e^-$.
    For the cases of $\pi^0$ and $\Sigma^0$, the slopes of the radiative corrections are denoted, from which one can infer the effects on the measured TFF slopes (see Eq.~\ref{eq:aQED}).
    }
    \label{fig:Dalitz_RC}
\end{figure}
%
This corresponds to the QED correction of about $-6\%$ to the slope estimate in absolute numbers.
This in turn means that without considering the QED corrections one would inevitably obtain a negative slope of similar size as the correct result $a_\pi$ stated above.

In QED, in which hadronic effects of the $\pi^0\to\gamma^*\gamma^*$ transition are included via the electromagnetic form factor, it turns out that one can precisely determine the Dalitz-decay branching ratio, directly following from the NLO calculation of the ratio \mbox{$R=B(\pi^0\to e^+e^-\gamma(\gamma))/B(\pi^0\to\gamma\gamma)$}~\cite{Husek:2018qdx}.
For this, it is sufficient to know the NLO radiative corrections and --- from theoretical estimates confirmed by experiment --- that the TFF slope is rather small.
This then combines with the fact that the sensitivity to the only relevant (higher derivatives can be neglected since the pion mass lies well below the resonance region) hadronic parameter present ($a_\pi$) in such an equation
%
\begin{multline}
    R
    =\frac\alpha\pi\int\!\!\!\!\int(1+a_\pi x)^2\frac{(1-x)^3}{4x}\\
    \times\big[1+\delta(x,y)\big]\biggl(1+y^2+\frac{4m_e^2}{M_\pi^2x}\biggr)\,\mathrm{d}x\mathrm{d}y\,
\end{multline}
%
is suppressed because the region with high virtualities is (and the slope thus does not influence the final result much); this is represented by the $(1-x)^3$ term above.
Let us stress that such a precise determination is only possible since the complete set of NLO QED radiative corrections are now known.

Believing the above determination, it becomes surprising that the KTeV Collaboration, soon after Ref.~\cite{Husek:2018qdx} was published, presented their results on the Dalitz-decay branching ratio~\cite{KTeV:2019gud} which is 3.6\,$\sigma$ away from the mentioned QED-based result.
It is thus desirable to have another independent branching-ratio measurement.
There is ongoing activity in this direction in NA62.


\subsubsection{\texorpdfstring{$\pi^0\to 2(e^+e^-)$}{pi0->4e}}

A complete set of NLO radiative corrections in soft-photon approximation for the neutral-pion double-Dalitz decay is computed in Ref.~\cite{Kampf:2018wau}, in which the processes $P\to\bar\ell\ell\bar\ell'\ell'$ are addressed in general.
This work also considers form-factor effects and improves upon the earlier work~\cite{Barker:2002ib}.


\subsection[\texorpdfstring{$\eta^{(\prime)}$}{eta(')} Dalitz decays]{Radiative corrections for \texorpdfstring{$\eta^{(\prime)}$}{eta(')} Dalitz decays}

The knowledge of the radiative corrections for the neutral-pion Dalitz decay turns out to be very useful in different contexts, and it also serves as a starting point to the calculation of the corrections for the corresponding $\eta^{(\prime)}$ decay channels.
These come with several additional challenges that one needs to overcome in order to be able to extract valuable information on the TFFs from data.

Naive radiative corrections for $\eta\to e^+e^-\gamma$~\cite{Mikaelian:1972jn} were indeed closely inspired by the solution for its neutral-pion counterpart, and only the mass of the decaying meson would be numerically addressed.
However, one needs to take the situation a bit more seriously.
The larger rest mass of the decaying $\eta^{(\prime)}$ has further consequences.
The $\eta$-meson mass already lies above the muon-pair and pion-pair production threshold, which has some consequences for the treatment of virtual corrections.
But what brings in a real complication is the fact that $\eta'$ lies above the lowest-lying resonances $\rho$ and $\omega$.
The decays and the radiative corrections then naturally become sensitive to the widths and shapes of the resonances.
And this is a crucial difference compared to the $\pi^0$ case, where the TFF dependence was mild and could even be neglected in certain applications.
Another technical complication then resides in the fact that $\eta^{(\prime)}$ mesons carry a nontrivial strangeness content that has to be addressed as part of the TFF treatment.

A more advanced approach to the $\eta^{(\prime)}\to\ell^+\ell^-\gamma$ radiative corrections~\cite{Husek:2017vmo} then improves on the naive approach for $\eta\to e^+e^-\gamma$ in the following way.
Regarding the vacuum polarization contribution, the muon-loop and hadronic contributions were added, the 1$\gamma$IR contribution at one-loop level was included, and the form-factor effects (including the complete bremsstrahlung contribution) were calculated.
Moreover, no assumptions regarding the final-state lepton masses were considered, so the muon channels could be consistently treated as well.
An example of the resulting radiative corrections can be seen in Fig.~\ref{fig:etap_e_RC}.
%
\begin{figure}[ht]
    \centering
    \includegraphics[width=\columnwidth]{Mesons/fig/etap_e_x.pdf}\\[2mm]
    \caption{Inclusive NLO QED radiative corrections for $\eta'\to e^+e^-\gamma$; \mbox{$x\equiv(p_{e^+}+p_{e^-})^2/M_{\eta'}^2$}.}
    \label{fig:etap_e_RC}
\end{figure}

Let us note that the developed exact treatment of the form-factor effects can also be used in the $\pi^0$ Dalitz decay.
There it introduces circa 1\% correction to the correction itself and is thus, as expected, negligible.


\subsection[\texorpdfstring{$\Sigma^0\to\Lambda e^+e^-$}{Sigma0->Lambda ee} Dalitz decay]{Radiative corrections for \texorpdfstring{$\Sigma^0\to\Lambda e^+e^-$}{Sigma0->Lambda ee} Dalitz decay}

The next technical complication in the calculation of radiative corrections to differential decay widths of Dalitz decays appears when the final-state photon is replaced by a massive particle.
The first example of such a case that we discuss here appears outside the meson sector: the Dalitz decay of the neutral Sigma baryon.

Radiative corrections within the soft-photon-approximation approach have been available for some time now~\cite{Sidhu:1971myo}.
The calculation of the bremsstrahlung contribution including the effects of hard-photon emission proceeds along similar lines as in the previous cases, although now the analytical integrals need to be generalized for a nonvanishing (Lambda-hyperon) mass.
It turns out that some of the integrals do not have analytical solutions anymore with such a generalization.
The counterpart of the (until now triangular) 1$\gamma$IR contribution of Fig.~\ref{fig:1gIR} are here two-photon-exchange box diagrams involving hadronic form factors; see Fig.~\ref{fig:1gIR_box}.
%
\begin{figure}[ht]
    \centering
    \begin{subfigure}[t]{0.48\columnwidth}
        \centering
        \includegraphics[width=\textwidth]{Mesons/fig/1gIR.eps}
    \end{subfigure}
    %
    \begin{subfigure}[t]{0.48\columnwidth}
        \centering
        \includegraphics[width=\textwidth]{Mesons/fig/1gIR_2.eps}
    \end{subfigure}\\[3mm]
    \caption{1$\gamma$IR (two-photon-exchange) Feynman diagrams contributing to NLO radiative corrections for~\mbox{$\Sigma^0\to\Lambda e^+e^-$}.}
    \label{fig:1gIR_box}
\end{figure}
%
Although quite nontrivial to evaluate and contributing to the electron--positron asymmetry, they do not affect the one-fold differential decay width in terms of the invariant electron--positron mass squared, and thus the extraction of the $\Sigma^0\to\Lambda\gamma^*$ transition form factor.
Here lies a profound difference between the kinematical regimes of decay and scattering since, in the latter, the two-photon-exchange contribution matters, as discussed in Sec.~\ref{sec:TPE}.

Similarly to the $\pi^0$-Dalitz-decay case, the (magnetic; electric is negligible) form-factor dependence is rather mild, and one can come up with precise predictions for $\Sigma^0\to\Lambda e^+e^-$ and $\Sigma^0\to\Lambda\gamma$ branching ratios~\cite{Husek:2019wmt}.
Regarding the importance of the QED NLO radiative corrections (see Fig.~\ref{fig:Dalitz_RC}), it was found, using again the simple consideration
%
\begin{equation}
    a_{(+\text{QED})}-a
    \simeq\frac12\frac{\mathrm{d}\delta(x)}{\mathrm{d}x}\bigg|_{x\ll1}\,,
\label{eq:aQED}
\end{equation}
%
that the effect on the TFF-slope extraction is about $-3.5\%$ in absolute figures~\cite{Husek:2019wmt}.
This is quite a significant value while extracting a quantity that is estimated to be of a size $1.8(3)$\%: When ignoring the QED effects, one would thus end up with a negative magnetic TFF slope, although again of a reasonable magnitude.


\subsection[\texorpdfstring{$\omega\to\pi^0 \ell^+\ell^-$}{omega->pi0ll}]{Radiative corrections for \texorpdfstring{$\omega\to\pi^0 \ell^+\ell^-$}{omega->pi0ll}}

Once the bremsstrahlung contribution to the $\Sigma^0$ Dalitz decay is established, the results can be readily used for the $\omega\to\pi^0 \ell^+\ell^-$ decay.
Here, various extractions of the $\omega\pi V$ correlator~\cite{Dzhelyadin:1980tj,NA60:2009una,NA60:2016nad,Adlarson:2016hpp} are in tension with the simplest theoretical models, and one could wonder if applying the previously omitted QED radiative corrections (they were not available until very recently) could improve the situation.
The corrections are, as in the previous cases, negative over a large region of the Dalitz plot (see Fig.~\ref{fig:Dalitz_RC}), and one can thus expect the data points to be pushed upwards accordingly after the QED effects are subtracted.
Needless to say, these corrections are more significant in the electron channel.
A proper analysis is, of course, in place and should be part of the future experimental analyses of more precise data.

The one-photon-inclusive NLO QED radiative corrections beyond the soft-photon approximation were calculated as a bachelor project in 2021~\cite{Lindahl:bc_thesis}.


\subsection[\texorpdfstring{$K^+\to\pi^+\ell^+\ell^-$}{K+->pi+ll}]{Radiative corrections for \texorpdfstring{$K^+\to\pi^+\ell^+\ell^-$}{K+->pi+ll}}

The $K^+\to\pi^+\ell^+\ell^-$ decays allow us to access the $K^+\to\pi^+\gamma^*$ TFF.
Extracted both in the electron and muon channels, it allows us to test the lepton-flavor universality.
Most importantly, however, studying the TFF provides us with valuable information about weak transitions in the low-energy QCD sector.
Extracting the relevant related hadronic parameters from experiment then requires, as expected, the knowledge of radiative corrections in order to subtract the QED effects that are present.
Recently, following a number of previous extractions~\cite{E865:1999ker,NA482:2009pfe,NA482:2010zrc}, the TFF parameters traditionally called $a_+$ and $b_+$ were measured by the NA62 Collaboration~\cite{NA62:2022qes}; it is worth mentioning that radiative corrections were part of this analysis at the MC level~\cite{Kubis:2010mp,Husek:2022fsc}.

Regarding the radiative corrections related to the lepton part of the $K^+\to\pi^+\ell^+\ell^-$ process, there are no additional complications, and one can use the same approach as in the previous cases once the form factor is suitably expanded to account for hadronic effects in the bremsstrahlung contribution {\em beyond} the soft-photon approximation.
On the other hand, the novel nontrivial obstacle is the fact that the involved kaon and pion are charged and can thus both also radiate a bremsstrahlung photon.
Moreover, the transition $K^+\to\pi^+\gamma^*\gamma^{(*)}$ needs to be considered.
Finally, virtual corrections that complement the bremsstrahlung from the meson part are somewhat nontrivial.
Everything then should be ideally expressed in terms of a single form factor, the one that is extracted from data.
This then allows for factorizing it out or for further iterations.
